\documentclass{veritas-epreuve}
\classe{3\textsuperscript{e}}
\matiere{Mathématiques}
\epreuve{Évaluation de la séquence 1}
\duree{2 heures}
\coefficient{4}
\code{VRT-3E-MATHS-S1}

\begin{document}
\titreepreuve

\begin{consignes}
\textbf{La calculatrice non programmable est autorisée.} Le sujet comporte trois exercices
indépendants, à traiter dans l'ordre de votre choix. La clarté des raisonnements et la
qualité de la rédaction entrent pour une part importante dans l'appréciation des copies.\par
\textbf{Barème} : exercice 1, 6 points ; exercice 2, 7 points ; exercice 3, 7 points.
\end{consignes}

\exercice[Calcul numérique]{6 points}
On donne $A = \dfrac{3}{4} - \dfrac{5}{6} \times \dfrac{9}{10}$ et
$B = \sqrt{75} - 2\sqrt{12} + \sqrt{27}$.
\begin{enumerate}
  \item Calculer $A$ et donner le résultat sous forme d'une fraction irréductible.
  \begin{solution}
    $A = \dfrac{3}{4} - \dfrac{45}{60} = \dfrac{3}{4} - \dfrac{3}{4} = 0$. \bareme{2 pts}
  \end{solution}
  \item Écrire $B$ sous la forme $a\sqrt{3}$, où $a$ est un entier relatif.
  \begin{solution}
    $\sqrt{75}=5\sqrt3$, $\sqrt{12}=2\sqrt3$, $\sqrt{27}=3\sqrt3$, donc
    $B = 5\sqrt3 - 4\sqrt3 + 3\sqrt3 = 4\sqrt3$. \bareme{2 pts}
  \end{solution}
  \item Montrer que $B^2$ est un nombre entier.
  \begin{solution}
    $B^2 = (4\sqrt3)^2 = 16 \times 3 = 48$, qui est entier. \bareme{2 pts}
  \end{solution}
\end{enumerate}

\exercice[Propriété de Thalès]{7 points}
\begin{minipage}[t]{.55\linewidth}
Sur la figure ci-contre, qui n'est pas en vraie grandeur, les points $A$, $M$, $B$ sont
alignés, ainsi que les points $A$, $N$, $C$. On donne :
$AM = 4$ cm, $AB = 10$ cm, $AN = 2{,}4$ cm, $AC = 6$ cm et $BC = 8$ cm.
\begin{enumerate}
  \item Démontrer que les droites $(MN)$ et $(BC)$ sont parallèles.
  \item Calculer la longueur $MN$.
  \item Le triangle $ABC$ est-il rectangle ? Justifier.
\end{enumerate}
\end{minipage}\hfill
\begin{minipage}[t]{.4\linewidth}
\vspace{-4mm}
\centering
\begin{tikzpicture}[figure,scale=.56]
  \coordinate (A) at (0,0);
  \coordinate (B) at (10,0);
  \coordinate (C) at (3.6,4.8);
  \coordinate (M) at ($(A)!.4!(B)$);
  \coordinate (N) at ($(A)!.4!(C)$);
  \draw (A)--(B)--(C)--cycle;
  \draw[vrtbleu] (M)--(N);
  \foreach \p/\pos in {A/below left,B/below right,C/above,M/below,N/above left}
    \node[point,label={\pos:$\p$}] at (\p) {};
\end{tikzpicture}
\end{minipage}

\begin{solution}
\begin{enumerate}
  \item On calcule
  \[ \frac{AM}{AB} = \frac{4}{10} = 0{,}4 \qquad\text{et}\qquad \frac{AN}{AC} = \frac{2{,}4}{6} = 0{,}4. \]
  Les points $A, M, B$ et $A, N, C$ sont alignés dans le même ordre et
  $AM/AB = AN/AC$ : d'après la réciproque de la propriété de Thalès,
  $(MN) \parallel (BC)$. \bareme{3 pts}
  \item D'après la propriété de Thalès, $\dfrac{MN}{BC} = 0{,}4$, donc
  $MN = 0{,}4 \times 8 = 3{,}2$ cm. \bareme{2 pts}
  \item Le plus grand côté est $[AB]$. $AC^2 + BC^2 = 36 + 64 = 100$ et $AB^2 = 10^2 = 100$.
  Comme $AB^2 = AC^2 + BC^2$, d'après la réciproque de la propriété de Pythagore, le triangle
  $ABC$ est rectangle en $C$. \bareme{2 pts}
\end{enumerate}
\end{solution}

\exercice[Fonction affine]{7 points}
Un taxi-moto de Douala propose deux tarifs pour un trajet de $x$ kilomètres :
\begin{itemize}
  \item \textbf{tarif A} : 150 F par kilomètre ;
  \item \textbf{tarif B} : une prise en charge de 300 F, puis 100 F par kilomètre.
\end{itemize}
\begin{enumerate}
  \item Exprimer, en fonction de $x$, le prix $f(x)$ au tarif A et le prix $g(x)$ au tarif B.
  \item Représenter $f$ et $g$ dans un même repère.
  \item Pour quelle distance les deux tarifs sont-ils égaux ? Lequel choisir pour 10 km ?
\end{enumerate}

\begin{solution}
\begin{enumerate}
  \item $f(x) = 150x$ et $g(x) = 100x + 300$. \bareme{2 pts}
  \item \hfill\begin{tikzpicture}[baseline=(current bounding box.north)]
    \begin{axis}[width=8.5cm,height=5.6cm,xmin=0,xmax=10.5,ymin=0,ymax=1650,
        xtick={0,2,...,10},ytick={0,300,...,1500},clip=false]
      \node[anchor=south east,font=\footnotesize] at ([yshift=3pt]rel axis cs:1,0) {$x$ (km)};
      \node[anchor=south west,font=\footnotesize] at (rel axis cs:0,1) {prix (F)};
      \addplot[vrtbleu,very thick,domain=0:10] {150*x} node[pos=.92,above left]{$f$};
      \addplot[vrtvert,very thick,domain=0:10] {100*x+300} node[pos=.92,below right]{$g$};
      \addplot[only marks,mark=*,mark size=1.6pt] coordinates {(6,900)};
      \draw[dashed,vrtcorps] (axis cs:6,0)--(axis cs:6,900)--(axis cs:0,900);
    \end{axis}
  \end{tikzpicture}\hfill\null \bareme{3 pts}
  \item $150x = 100x + 300 \iff 50x = 300 \iff x = 6$ : les tarifs sont égaux pour 6 km
  (900 F). Pour 10 km : $f(10) = 1\,500$ F et $g(10) = 1\,300$ F ; le tarif B est le
  plus avantageux. \bareme{2 pts}
\end{enumerate}
\end{solution}

\end{document}
